% GENERATED FILE. Edit canonical lessons, then run npm run generate:books.
% canonical-source-hash: fnv1a64:ecd4ebfe4b0c7680
% canonical-lessons: KA-C540-kabaddi, KA-C540-kusti, KA-C540-golu, KA-C540-ode-kick, KA-C540-guddu

\chapter{Kicks and Punches}
\label{ch:ka-kicks-and-punches}
\hlchaptermodality{540}

\begin{chapteropening}
\textbf{By the end of this chapter:} \emph{I can say kabaddi, wrestling, a goal in sport, to kick and to punch.}

Complete the last lesson of chapter 540: I can say kabaddi, wrestling, a goal in sport, to kick and to punch.
\end{chapteropening}

\section[\texorpdfstring{\kn{ಕಬಡ್ಡಿ} (kabaḍḍi)}{ಕಬಡ್ಡಿ (kabaḍḍi)}]{\kn{ಕಬಡ್ಡಿ} (kabaḍḍi) — kabaddi}

\label{lesson:KA-C540-kabaddi}



\begin{quote}
Before the new one: say the Kannada for a lead, then the Kannada for to contend.
\end{quote}

\subsection*{You'll want to know: \kn{ಕಬಡ್ಡಿ}}
\textbf{\kn{ಕಬಡ್ಡಿ}} — \emph{kabaḍḍi} — \textquotedblleft{}kabaddi\textquotedblright{}.

\begin{quote}
\kn{ಹಳ್ಳಿಯ} \kn{ಹುಡುಗರು} \kn{ಮೈದಾನದಲ್ಲಿ} \kn{ಕಬಡ್ಡಿ} \kn{ಆಡುತ್ತಾರೆ}.
\end{quote}

\begin{quote}
\emph{haḷḷiya huḍugaru maidānadalli kabaḍḍi āḍuttāre.}
\end{quote}

\begin{quote}
The village boys play kabaddi on the ground.
\end{quote}

A raider holds his breath, chants, and tags players on the other side.

\begin{cousinweb}
\textbf{\kn{ಕಬಡ್ಡಿ}} (\emph{kabaḍḍi}) is the name of a team game of south Asia; the origin of the name is not settled.
\end{cousinweb}

\begin{grammarlens}[title={What you've built}]
One more word for describing what you lived and felt.
\end{grammarlens}

\subsection*{Your turn}
\begin{itemize}
\raggedright
  \item \emph{Say it:} \emph{kabaḍḍi}
  \item \emph{Say it:} \emph{kabaḍḍi}, once more
  \item \emph{Say it:} \emph{haḷḷiya huḍugaru maidānadalli kabaḍḍi āḍuttāre.} — The village boys play kabaddi on the ground.
  \item \emph{Recall:} say the Kannada for a lead, then the Kannada for to contend, then say \emph{kabaḍḍi} again
\end{itemize}

\begin{tcolorbox}[breakable,colback=teal!4,colframe=teal!35!black,title={Before you move on}]
What does \kn{ಕಬಡ್ಡಿ} mean? (\textquotedblleft{}Kabaddi\textquotedblright{}.) Say it once more.
\end{tcolorbox}
\section[\texorpdfstring{\kn{ಕುಸ್ತಿ} (kusti)}{ಕುಸ್ತಿ (kusti)}]{\kn{ಕುಸ್ತಿ} (kusti) — wrestling}

\label{lesson:KA-C540-kusti}



\begin{quote}
Before the new one: say the Kannada for to contend, then the Kannada for kabaddi.
\end{quote}

\subsection*{You'll want to know: \kn{ಕುಸ್ತಿ}}
\textbf{\kn{ಕುಸ್ತಿ}} — \emph{kusti} — \textquotedblleft{}wrestling\textquotedblright{}.

\begin{quote}
\kn{ಜಾತ್ರೆಯಲ್ಲಿ} \kn{ದೊಡ್ಡ} \kn{ಕುಸ್ತಿ} \kn{ಸ್ಪರ್ಧೆ} \kn{ನಡೆಯಿತು}.
\end{quote}

\begin{quote}
\emph{jātreyalli doḍḍa kusti spardhe naḍeyitu.}
\end{quote}

\begin{quote}
A big wrestling contest was held at the fair.
\end{quote}

\kn{ಕುಸ್ತಿಪಟು}, a wrestler.

\begin{cousinweb}
\textbf{\kn{ಕುಸ್ತಿ}} (\emph{kusti}) is usually traced to Persian \emph{kushtī}, \textquotedblleft{}wrestling\textquotedblright{}, which is thought to have reached Kannada through Hindustani.
\end{cousinweb}

\begin{grammarlens}[title={What you've built}]
One more word for describing what you lived and felt.
\end{grammarlens}

\subsection*{Your turn}
\begin{itemize}
\raggedright
  \item \emph{Say it:} \emph{kusti}
  \item \emph{Say it:} \emph{kusti}, once more
  \item \emph{Say it:} \emph{jātreyalli doḍḍa kusti spardhe naḍeyitu.} — A big wrestling contest was held at the fair.
  \item \emph{Recall:} say the Kannada for to contend, then the Kannada for kabaddi, then say \emph{kusti} again
\end{itemize}

\begin{tcolorbox}[breakable,colback=teal!4,colframe=teal!35!black,title={Before you move on}]
What does \kn{ಕುಸ್ತಿ} mean? (\textquotedblleft{}Wrestling\textquotedblright{}.) Say it once more.
\end{tcolorbox}
\section[\texorpdfstring{\kn{ಗೋಲು} (gōlu)}{ಗೋಲು (gōlu)}]{\kn{ಗೋಲು} (gōlu) — a goal in sport}

\label{lesson:KA-C540-golu}



\begin{quote}
Before the new one: say the Kannada for kabaddi, then the Kannada for wrestling.
\end{quote}

\subsection*{You'll want to know: \kn{ಗೋಲು}}
\textbf{\kn{ಗೋಲು}} — \emph{gōlu} — \textquotedblleft{}a goal in sport\textquotedblright{}.

\begin{quote}
\kn{ಕೊನೆಯ} \kn{ನಿಮಿಷದಲ್ಲಿ} \kn{ಅವನು} \kn{ಗೋಲು} \kn{ಹೊಡೆದನು}.
\end{quote}

\begin{quote}
\emph{koneya nimiṣadalli avanu gōlu hoḍedanu.}
\end{quote}

\begin{quote}
He scored a goal in the last minute.
\end{quote}

\kn{ಗೋಲು} \kn{ಹೊಡೆ}, to score a goal. A goal in life is \kn{ಗುರಿ}.

\begin{cousinweb}
\textbf{\kn{ಗೋಲು}} (\emph{gōlu}) is a loan from English \emph{goal}, with the -\kn{ು} many loans take in speech.
\end{cousinweb}

\begin{grammarlens}[title={What you've built}]
One more word for describing what you lived and felt.
\end{grammarlens}

\subsection*{Your turn}
\begin{itemize}
\raggedright
  \item \emph{Say it:} \emph{gōlu}
  \item \emph{Say it:} \emph{gōlu}, once more
  \item \emph{Say it:} \emph{koneya nimiṣadalli avanu gōlu hoḍedanu.} — He scored a goal in the last minute.
  \item \emph{Recall:} say the Kannada for kabaddi, then the Kannada for wrestling, then say \emph{gōlu} again
\end{itemize}

\begin{tcolorbox}[breakable,colback=teal!4,colframe=teal!35!black,title={Before you move on}]
What does \kn{ಗೋಲು} mean? (\textquotedblleft{}A goal in sport\textquotedblright{}.) Say it once more.
\end{tcolorbox}
\section[\texorpdfstring{\kn{ಒದೆ} (ode)}{ಒದೆ (ode)}]{\kn{ಒದೆ} (ode) — to kick}

\label{lesson:KA-C540-ode-kick}



\begin{quote}
Before the new one: say the Kannada for wrestling, then the Kannada for a goal in sport.
\end{quote}

\subsection*{You'll want to know: \kn{ಒದೆ}}
\textbf{\kn{ಒದೆ}} — \emph{ode} — \textquotedblleft{}to kick\textquotedblright{}.

\begin{quote}
\kn{ಹುಡುಗ} \kn{ಚೆಂಡನ್ನು} \kn{ಜೋರಾಗಿ} \kn{ಒದ್ದನು}.
\end{quote}

\begin{quote}
\emph{huḍuga ceṇḍannu jōrāgi oddanu.}
\end{quote}

\begin{quote}
The boy kicked the ball hard.
\end{quote}

Past (I …): \kn{ಒದ್ದೆ}, \emph{odde}.

The past \kn{ಒದ್ದೆ}, I kicked, is spelled like \kn{ಒದ್ದೆ}, wet; the sentence tells them apart.

\begin{cousinweb}
\textbf{\kn{ಒದೆ}} (\emph{ode}) is native Dravidian: DEDR 616 lists Kannada \emph{ode}, \emph{odi}, \textquotedblleft{}to kick, spur\textquotedblright{}, with \emph{odd-} in the past.
\end{cousinweb}

\begin{grammarlens}[title={What you've built}]
One more word for describing what you lived and felt.
\end{grammarlens}

\subsection*{Your turn}
\begin{itemize}
\raggedright
  \item \emph{Say it:} \emph{ode}
  \item \emph{Say it:} \emph{ode}, once more
  \item \emph{Say it:} \emph{huḍuga ceṇḍannu jōrāgi oddanu.} — The boy kicked the ball hard.
  \item \emph{Recall:} say the Kannada for wrestling, then the Kannada for a goal in sport, then say \emph{ode} again
\end{itemize}

\begin{tcolorbox}[breakable,colback=teal!4,colframe=teal!35!black,title={Before you move on}]
What does \kn{ಒದೆ} mean? (\textquotedblleft{}To kick\textquotedblright{}.) Say it once more.
\end{tcolorbox}
\section[\texorpdfstring{\kn{ಗುದ್ದು} (guddu)}{ಗುದ್ದು (guddu)}]{\kn{ಗುದ್ದು} (guddu) — to punch}

\label{lesson:KA-C540-guddu}



\begin{quote}
Before the new one: say the Kannada for a goal in sport, then the Kannada for to kick.
\end{quote}

\subsection*{You'll want to know: \kn{ಗುದ್ದು}}
\textbf{\kn{ಗುದ್ದು}} — \emph{guddu} — \textquotedblleft{}to punch\textquotedblright{}.

\begin{quote}
\kn{ಆಟದಲ್ಲಿ} \kn{ತಮ್ಮ} \kn{ನನ್ನ} \kn{ಬೆನ್ನಿಗೆ} \kn{ಗುದ್ದಿದನು}.
\end{quote}

\begin{quote}
\emph{āṭadalli tamma nanna bennige guddidanu.}
\end{quote}

\begin{quote}
In the game my younger brother punched my back.
\end{quote}

Past (I …): \kn{ಗುದ್ದಿದೆ}, \emph{guddide}.

Also a noun, a punch.

\begin{cousinweb}
\textbf{\kn{ಗುದ್ದು}} (\emph{guddu}) is native Dravidian: DEDR 1850 lists Kannada \emph{kuttu}, \textquotedblleft{}to beat, strike\textquotedblright{}, and \emph{guddu}, \textquotedblleft{}to strike with the fist\textquotedblright{}.
\end{cousinweb}

\begin{grammarlens}[title={What you've built}]
One more word for describing what you lived and felt.
\end{grammarlens}

\subsection*{Your turn}
\begin{itemize}
\raggedright
  \item \emph{Say it:} \emph{guddu}
  \item \emph{Say it:} \emph{guddu}, once more
  \item \emph{Say it:} \emph{āṭadalli tamma nanna bennige guddidanu.} — In the game my younger brother punched my back.
  \item \emph{Recall:} say the Kannada for a goal in sport, then the Kannada for to kick, then say \emph{guddu} again
\end{itemize}

\begin{tcolorbox}[breakable,colback=teal!4,colframe=teal!35!black,title={Before you move on}]
What does \kn{ಗುದ್ದು} mean? (\textquotedblleft{}To punch\textquotedblright{}.) Say it once more.
\end{tcolorbox}
